\chapter{ทฤษฎีสัมพัทธภาพพิเศษ}
\label{chap:special_relativity}
\index{สัมพัทธภาพพิเศษ}
\index{ไอน์สไตน์}

% ==========================================================
% แผนบริหารการสอนประจำบทที่ 2
% ==========================================================
\section*{แผนบริหารการสอนประจำบทที่ 2}
\addcontentsline{toc}{section}{แผนบริหารการสอนประจำบทที่ 2}

\begin{lessonplanbox}
\noindent\textbf{หัวข้อเนื้อหาประจำบท}
\begin{enumerate}
    \item วิกฤตการณ์ฟิสิกส์ดั้งเดิมและการทดลองของไมเคิลสัน-มอร์ลีย์
    \item สัจพจน์ของไอน์สไตน์
    \item การแปลงแบบลอเรนซ์
    \item ผลสืบเนื่องทางจลนศาสตร์ การยืดออกของเวลาและการหดสั้นของความยาว
    \item ปริภูมิ-เวลา 4 มิติ และกรวยแสง
    \item พลศาสตร์สัมพัทธภาพ โมเมนตัมและพลังงานรวม
    \item สถาปัตยกรรมการจำลองเสมือนจริง AR/XR 4D Spacetime Visualizer
    \item การคำนวณเชิงตัวเลขด้วย Python การวิเคราะห์ตัวคูณลอเรนซ์
    \item ตัวอย่างการแก้โจทย์เชิงวิเคราะห์ตามระเบียบวิธี P-S-C-T
    \item สรุปสาระสำคัญและแบบฝึกหัดท้ายบท
\end{enumerate}

\vspace{0.4em}
\noindent\textbf{วัตถุประสงค์เชิงพฤติกรรม}
\begin{enumerate}
    \item อธิบายข้อจำกัดของกลศาสตร์นิวตันและการทดลองของไมเคิลสัน-มอร์ลีย์ได้อย่างถูกต้อง
    \item สรุปสัจพจน์ของไอน์สไตน์ทั้งสองประการ และพิสูจน์ที่มาของการแปลงแบบลอเรนซ์ได้อย่างเป็นระบบ
    \item คำนวณปรากฏการณ์การยืดออกของเวลาและการหดสั้นของความยาวได้อย่างแม่นยำ
    \item วิเคราะห์โครงสร้างปริภูมิ-เวลา 4 มิติและแผนภาพกรวยแสงได้อย่างถูกต้อง
    \item คำนวณความสัมพันธ์ระหว่างโมเมนตัมสัมพัทธภาพและพลังงานรวม $E^2 = p^2 c^2 + m_0^2 c^4$ ได้อย่างถูกต้อง
    \item ออกแบบและพัฒนาสคริปต์ Python เพื่อจำลองการแปลงพิกัดและการพล็อตกราฟตัวคูณลอเรนซ์ $\gamma(v)$ ได้อย่างมีประสิทธิภาพ
\end{enumerate}

\vspace{0.4em}
\noindent\textbf{กิจกรรมการเรียนการสอน}
\begin{enumerate}
    \item การบรรยายเชิงทฤษฎีควบคู่กับการอภิปรายผลการทดลองของไมเคิลสัน-มอร์ลีย์
    \item กิจกรรมกลุ่มย่อยวิเคราะห์การแปลงลอเรนซ์และการยืดออกของเวลาผ่านการทดลองทางความคิดนาฬิกาแสง
    \item ปฏิบัติการจำลองสถานการณ์เสมือนจริง 4D Spacetime Visualizer บนระบบ WebXR
    \item ปฏิบัติการเขียนโค้ดภาษา Python จำลองการหดสั้นของความยาวและตัวคูณลอเรนซ์
    \item การศึกษาและแก้โจทย์ปัญหาฟิสิกส์สัมพัทธภาพตามระเบียบวิธี P-S-C-T
\end{enumerate}

\vspace{0.4em}
\noindent\textbf{สื่อการเรียนการสอน}
\begin{enumerate}
    \item เอกสารคำสอน บทที่ 2 ทฤษฎีสัมพัทธภาพพิเศษ
    \item สื่อจำลอง 3 มิติปฏิสัมพันธ์เสมือนจริง WebXR 4D Spacetime Visualizer
    \item ชุดสคริปต์ภาษา Python คำนวณและพล็อตกราฟพลศาสตร์สัมพัทธภาพ
    \item แพลตฟอร์มการจัดการเรียนรู้ดิจิทัล (LMS) และชุดโจทย์ปัญหา P-S-C-T
\end{enumerate}

\vspace{0.4em}
\noindent\textbf{การประเมินผล}
\begin{enumerate}
    \item การประเมินจากการตอบคำถามและการมีส่วนร่วมในกิจกรรมการเรียนรู้ในชั้นเรียน
    \item การตรวจรายงานผลปฏิบัติการคอมพิวเตอร์ด้วย Python และแบบฝึกหัดท้ายบท
    \item การประเมินโครงงานจำลองสถานการณ์ 3 มิติเสมือนจริง WebXR
    \item การทดสอบย่อยวัดผลสัมฤทธิ์ทางการเรียนประจำบทที่ 2
\end{enumerate}
\end{lessonplanbox}

\newpage

\section{วิกฤตการณ์ฟิสิกส์ดั้งเดิมและการทดลองของไมเคิลสัน-มอร์ลีย์}
\index{ไมเคิลสัน-มอร์ลีย์}
\index{อีเทอร์}

ในช่วงปลายคริสต์ศตวรรษที่ 19 ทฤษฎีแม่เหล็กไฟฟ้าของเจมส์ คลาร์ก แมกซ์เวลล์ (James Clerk Maxwell) ระบุว่า คลื่นแม่เหล็กไฟฟ้าเคลื่อนที่ด้วยอัตราเร็วในสุญญากาศคงที่เสมอคือ
\begin{equation}
c = \frac{1}{\sqrt{\mu_0 \varepsilon_0}} \approx 2.99792458 \times 10^8 \text{ m/s}
\end{equation}
ทว่าในกรอบกลศาสตร์ดั้งเดิมของกาลิเลโอและนิวตัน ความเร็วใดๆ ย่อมขึ้นอยู่กับกรอบอ้างอิงเฉื่อย (Inertial Reference Frame) ซึ่งก่อให้เกิดสมมติฐานเกี่ยวกับตัวกลางพาแสงที่เรียกว่า \textbf{อีเทอร์เรืองแสง (Luminiferous Aether)}

อัลเบิร์ต ไมเคิลสัน (Albert A. Michelson) และเอ็ดเวิร์ด มอร์ลีย์ (Edward W. Morley) ได้ทำการทดลองอันลือลั่นในปี ค.ศ. 1887 โดยใช้เครื่องแทรกสอดแสง (Interferometer) เพื่อตรวจจับ "ลมของอีเทอร์" (Aether Drift) จากการเคลื่อนที่ของโลกรอบดวงอาทิตย์ ผลลัพธ์ที่ได้ปรากฏเป็น \textit{ศูนย์ (Null Result)} อย่างน่าทึ่ง นั่นคือแสงมีความเร็วเท่าเดิมในทุกทิศทาง โดยไม่ขึ้นกับการเคลื่อนที่ของโลกเลยแม้แต่น้อย

\begin{figure}[H]
\centering
\begin{tikzpicture}[scale=0.95]
    % Source and beam splitter
    \draw[darkNavy, thick, fill=cyan!15] (0,0) rectangle (0.8,0.8);
    \node at (0.4, 0.4) [font=\footnotesize\bfseries] {แหล่งกำเนิด};
    \draw[-{Stealth[scale=1.2]}, red!80!black, line width=1.5pt] (0.8, 0.4) -- (3.0, 0.4);
    
    % Beam splitter (semi-transparent mirror)
    \draw[thick, fill=blue!10, rotate around={45:(3.5,0.4)}] (3.2, -0.3) rectangle (3.8, 1.1);
    \node at (3.5, -0.3) [font=\scriptsize, align=center] {แผ่นแบ่งลำแสง\\(Beam Splitter)};
    
    % Transmitted path to Mirror 1
    \draw[-{Stealth[scale=1.2]}, red!80!black, line width=1.2pt] (3.5, 0.4) -- (6.5, 0.4);
    \draw[{Stealth[scale=1.2]}-, blue!80!black, line width=1.2pt, dashed] (3.5, 0.5) -- (6.5, 0.5);
    \draw[darkNavy, fill=gray!40, thick] (6.5, -0.2) rectangle (6.7, 1.0);
    \node at (6.6, 1.3) [font=\footnotesize\bfseries] {กระจก $M_1$};
    
    % Reflected path to Mirror 2
    \draw[-{Stealth[scale=1.2]}, red!80!black, line width=1.2pt] (3.5, 0.4) -- (3.5, 3.4);
    \draw[{Stealth[scale=1.2]}-, blue!80!black, line width=1.2pt, dashed] (3.6, 0.4) -- (3.6, 3.4);
    \draw[darkNavy, fill=gray!40, thick] (2.9, 3.4) rectangle (4.1, 3.6);
    \node at (3.5, 3.9) [font=\footnotesize\bfseries] {กระจก $M_2$};
    
    % Combined beam to Detector
    \draw[-{Stealth[scale=1.2]}, purple, line width=1.5pt] (3.5, 0.4) -- (3.5, -2.0);
    \draw[darkNavy, fill=quantumGold!30, thick] (2.8, -2.4) rectangle (4.2, -2.0);
    \node at (3.5, -2.2) [font=\footnotesize\bfseries] {เครื่องตรวจวัด (Detector)};
    
    % Distances L1 and L2
    \draw[<->, slateGrey, thick] (3.5, 0.9) -- (6.5, 0.9) node[midway, above, font=\scriptsize] {$L_1$};
    \draw[<->, slateGrey, thick] (4.0, 0.4) -- (4.0, 3.4) node[midway, right, font=\scriptsize] {$L_2$};
\end{tikzpicture}
\caption{ผังการทดลองของไมเคิลสัน-มอร์ลีย์ (Michelson-Morley Experiment) ซึ่งแสดงผลลัพธ์เป็นศูนย์}
\label{fig:michelson_morley}
\end{figure}

\section{สัจพจน์ของไอน์สไตน์}
\index{สัจพจน์ของไอน์สไตน์}

ในปี ค.ศ. 1905 อัลเบิร์ต ไอน์สไตน์ ได้เสนอแนวคิดปฏิวัติวงการฟิสิกส์ด้วย \textbf{สัจพจน์ของไอน์สไตน์ (Einstein's Postulates)} โดยตั้งสมมติฐานมูลฐาน 2 ข้อ

\begin{mathdefinition}[สัจพจน์สองประการของทฤษฎีสัมพัทธภาพพิเศษ]
\begin{enumerate}
    \item \textbf{หลักความสัมพัทธ์ (Principle of Relativity)}\quad กฎเกณฑ์ทางฟิสิกส์ทั้งปวงย่อมมีรูปแบบทางคณิตศาสตร์เหมือนกันในทุกกรอบอ้างอิงเฉื่อย (Inertial Frames of Reference) ไม่มีกรอบอ้างอิงใดที่มีอภิสิทธิ์เหนือกรอบอื่น
    \item \textbf{ความคงที่ของอัตราเร็วแสง (Constancy of the Speed of Light)}\quad อัตราเร็วของแสงในสุญญากาศ $c$ มีค่าคงที่สัมบูรณ์เท่ากันเสมอในทุกกรอบอ้างอิงเฉื่อย โดยไม่ขึ้นกับสถานะการเคลื่อนที่ของแหล่งกำเนิดแสงหรือผู้สังเกต
\end{enumerate}
\end{mathdefinition}

\section{การแปลงแบบลอเรนซ์}
\index{การแปลงแบบลอเรนซ์}
\index{ตัวคูณลอเรนซ์}

พิจารณากรอบอ้างอิงสองกรอบ $S$ และ $S'$ โดยที่ $S'$ เคลื่อนที่สัมพันธ์กับ $S$ ด้วยความเร็วคงที่ $v$ ในแนวแกน $+x$ โดยสมมติว่าจุดกำเนิด $O$ และ $O'$ ทับกันพอดีที่เวลา $t = t' = 0$ การแปลงพิกัดของกาลิเลโอดั้งเดิม ($x' = x - vt, t' = t$) ไม่สามารถรักษาความคงที่ของ $c$ ได้ ไอน์สไตน์จึงอนุมาน \textbf{การแปลงแบบลอเรนซ์ (Lorentz Transformations)} ที่ถูกต้อง

\begin{maththeorem}[สมการการแปลงแบบลอเรนซ์]
การแปลงพิกัดและเวลาระหว่างกรอบ $S(x, y, z, t)$ และกรอบ $S'(x', y', z', t')$ ถูกกำหนดโดย
\begin{align}
x' &= \gamma (x - vt) \\
y' &= y \\
z' &= z \\
t' &= \gamma \left( t - \frac{vx}{c^2} \right)
\end{align}
โดยที่ $\gamma$ คือ \textbf{ตัวคูณลอเรนซ์ (Lorentz Factor)}
\begin{equation}
\gamma = \frac{1}{\sqrt{1 - \beta^2}} = \frac{1}{\sqrt{1 - \frac{v^2}{c^2}}}, \quad \beta = \frac{v}{c}
\end{equation}
และการแปลงผกผัน (Inverse Lorentz Transformation) จาก $S'$ กลับสู่ $S$ ทำได้โดยเปลี่ยนเครื่องหมายของ $v$
\begin{align}
x &= \gamma (x' + vt') \\
t &= \gamma \left( t' + \frac{vx'}{c^2} \right)
\end{align}
\end{maththeorem}

\section[ผลสืบเนื่องทางจลนศาสตร์ การยืดออกของเวลาและการหดสั้นของความยาว]{ผลสืบเนื่องทางจลนศาสตร์ การยืดออกของเวลา\\และการหดสั้นของความยาว}
\index{การยืดออกของเวลา}
\index{การหดสั้นของความยาว}
\index{เวลาเฉพาะ}

\subsection{การยืดออกของเวลา}
ปรากฏการณ์\textbf{การยืดออกของเวลา (Time Dilation)} พิจารณานาฬิกาแสงที่อยู่นิ่งในกรอบ $S'$ ช่วงเวลาที่นาฬิกาบันทึกได้ในกรอบที่ตัวเองอยู่นิ่งเรียกว่า \textbf{เวลาเฉพาะ (Proper Time, $\Delta t_0$)} เมื่อผู้สังเกตในกรอบ $S$ วัดช่วงเวลาของเหตุการณ์เดียวกัน จะพบว่าเวลาเดินช้าลงตามความสัมพันธ์
\begin{equation}
\Delta t = \gamma \Delta t_0 = \frac{\Delta t_0}{\sqrt{1 - v^2/c^2}} \ge \Delta t_0
\end{equation}

\begin{figure}[H]
\centering
\begin{tikzpicture}[scale=0.95]
    % Left: Rest Frame S'
    \begin{scope}[shift={(-4.5, 0)}]
        \node at (0, 3.8) [font=\footnotesize\bfseries\color{darkNavy}] {(ก) กรอบอยู่นิ่ง $S'$ (เวลาเฉพาะ $\Delta t_0$)};
        % Mirrors
        \draw[darkNavy, fill=gray!30, thick] (-1.0, 3.0) rectangle (1.0, 3.2);
        \draw[darkNavy, fill=gray!30, thick] (-1.0, 0.0) rectangle (1.0, 0.2);
        \node at (0, 3.4) [font=\scriptsize] {กระจกบน};
        \node at (0, -0.2) [font=\scriptsize] {กระจกล่าง / แหล่งกำเนิด};
        % Light path up and down
        \draw[-{Stealth[scale=1.1]}, line width=1.4pt, color=quantumGold] (-0.15, 0.2) -- (-0.15, 3.0);
        \draw[-{Stealth[scale=1.1]}, line width=1.4pt, color=quantumGold, dashed] (0.15, 3.0) -- (0.15, 0.2);
        \draw[<->, slateGrey, thick] (-1.4, 0.2) -- (-1.4, 3.0) node[midway, left, font=\scriptsize] {$D$};
        \node at (0, -0.8) [font=\footnotesize] {$\Delta t_0 = \dfrac{2D}{c}$};
    \end{scope}

    % Divider
    \draw[slateGrey!40, dashed, thick] (-1.0, -1.0) -- (-1.0, 4.0);

    % Right: Moving Frame S
    \begin{scope}[shift={(1.0, 0)}]
        \node at (2.5, 3.8) [font=\footnotesize\bfseries\color{darkNavy}] {(ข) กรอบเคลื่อนที่ $S$ (เวลา $\Delta t = \gamma \Delta t_0$)};
        % Mirrors at t=0, t=dt/2, t=dt
        \draw[gray!50, fill=gray!15, dashed] (-0.8, 0.0) rectangle (0.8, 0.2);
        \draw[darkNavy, fill=gray!30, thick] (1.7, 3.0) rectangle (3.3, 3.2);
        \draw[gray!50, fill=gray!15, dashed] (4.2, 0.0) rectangle (5.8, 0.2);
        
        % Triangular light path
        \draw[-{Stealth[scale=1.2]}, line width=1.5pt, color=cyberBlue] (0, 0.2) -- (2.5, 3.0);
        \draw[-{Stealth[scale=1.2]}, line width=1.5pt, color=cyberBlue, dashed] (2.5, 3.0) -- (5.0, 0.2);
        
        % Geometry labels
        \node at (1.0, 1.8) [font=\scriptsize\color{cyberBlue}, rotate=48] {$L = c\Delta t / 2$};
        \draw[<->, slateGrey, thick] (2.5, 0.2) -- (2.5, 3.0) node[midway, right, font=\scriptsize] {$D$};
        \draw[<->, slateGrey, thick] (0, -0.3) -- (2.5, -0.3) node[midway, below, font=\scriptsize] {$v\Delta t / 2$};
        \draw[<->, slateGrey, thick] (0, -0.8) -- (5.0, -0.8) node[midway, below, font=\scriptsize] {$v\Delta t$};
        \node at (2.5, -1.5) [font=\footnotesize] {$c^2 (\Delta t/2)^2 = D^2 + v^2 (\Delta t/2)^2$};
    \end{scope}
\end{tikzpicture}
\caption{การทดลองทางความคิดเรื่องนาฬิกาแสง (Light Clock) พิสูจน์ปรากฏการณ์การยืดออกของเวลา}
\label{fig:light_clock}
\end{figure}

\subsection{การหดสั้นของความยาว}
ปรากฏการณ์\textbf{การหดสั้นของความยาว (Length Contraction)} วัตถุที่มีความยาววัดได้ขณะอยู่นิ่งในกรอบของตัวเองเรียกว่า \textbf{ความยาวเฉพาะ (Proper Length, $L_0$)} เมื่อผู้สังเกตที่เคลื่อนที่สัมพันธ์กับวัตถุด้วยความเร็ว $v$ ในแนวขนานกับมิติของวัตถุวัดความยาว จะพบว่าความยาวหดสั้นลง
\begin{equation}
L = \frac{L_0}{\gamma} = L_0 \sqrt{1 - \frac{v^2}{c^2}} \le L_0
\end{equation}
โปรดสังเกตว่า มิติในแนวตั้งฉากกับการเคลื่อนที่ ($y$ และ $z$) จะไม่มีการเปลี่ยนแปลงความยาวเลย ($L_\perp = L_{0\perp}$)

\begin{pitfallbox}[ความเข้าใจผิดเกี่ยวกับความยาวเฉพาะและเวลาเฉพาะ]
ผู้เริ่มต้นศึกษามักสับสนว่า "ใครเห็นเวลาช้าลง" หรือ "สิ่งของหดสั้นจริงหรือไม่ทางกายภาพ"
\begin{itemize}
    \item เวลาเฉพาะ $\Delta t_0$ ต้องเป็นเวลาที่วัดจาก \textbf{นาฬิกาเรือนเดียวที่อยู่ประจำที่ตำแหน่งเดิมของเหตุการณ์ทั้งสอง} เสมอ
    \item ความยาวเฉพาะ $L_0$ ต้องวัดในกรอบอ้างอิงที่วัตถุนั้นอยู่นิ่งสัมพัทธ์
    \item การหดสั้นไม่ใช่การบีบอัดทางกายภาพของเนื้อสาร แต่เกิดจากคุณสมบัติทางเรขาคณิตของปริภูมิ-เวลา (Spacetime) เมื่อสังเกตจากต่างกรอบอ้างอิง
\end{itemize}
\end{pitfallbox}

\section{ปริภูมิ-เวลา 4 มิติ และกรวยแสง}
\index{มินคอฟสกี}
\index{กรวยแสง}
\index{ช่วงปริภูมิเวลา}

แฮร์มันน์ มินคอฟสกี (Hermann Minkowski) ได้เสนอการรวมปริภูมิสามมิติเข้ากับเวลาเป็น \textbf{กาลอวกาศ 4 มิติ (Spacetime Continuum)} โดยกำหนดพิกัด 4 มิติเป็น $x^\mu = (ct, x, y, z)$ ช่วงปริภูมิเวลาสัมบูรณ์ (Invariant Spacetime Interval) ระหว่างสองเหตุการณ์นิยามโดย
\begin{equation}
\Delta s^2 = c^2 \Delta t^2 - \left( \Delta x^2 + \Delta y^2 + \Delta z^2 \right)
\end{equation}
ค่า $\Delta s^2$ มีคุณสมบัติคงตัวแบบลอเรนซ์ (Lorentz Invariant) ซึ่งแบ่งออกเป็น 3 ย่าน
\begin{enumerate}
    \item $\Delta s^2 > 0$\quad \textbf{เหมือนเวลา (Timelike)} เหตุการณ์ทั้งสองสามารถเชื่อมโยงกันด้วยความสัมพันธ์เชิงเหตุและผล (Causality) ได้
    \item $\Delta s^2 = 0$\quad \textbf{เหมือนแสง (Lightlike หรือ Null)} เส้นทางของอนุภาคที่เคลื่อนที่ด้วยความเร็วแสง (เช่น โฟตอน)
    \item $\Delta s^2 < 0$\quad \textbf{เหมือนปริภูมิ (Spacelike)} ไม่มีสัญญาณใดในเอกภพสามารถส่งหากันได้ทัน (ไม่มีความสัมพันธ์เชิงสาเหตุ)
\end{enumerate}

\begin{figure}[H]
\centering
\begin{tikzpicture}[scale=1.1]
    % Axes
    \draw[-{Stealth[scale=1.2]}, slateGrey, thick] (-3.5,0) -- (3.5,0) node[right, font=\footnotesize\bfseries] {$x$ (Space)};
    \draw[-{Stealth[scale=1.2]}, slateGrey, thick] (0,-2.5) -- (0,3.5) node[above, font=\footnotesize\bfseries] {$ct$ (Time)};
    
    % Light cone lines
    \draw[cyan, line width=1.5pt] (-2.8,-2.8) -- (2.8,2.8);
    \draw[cyan, line width=1.5pt] (-2.8,2.8) -- (2.8,-2.8);
    
    % Shading Future Light Cone
    \fill[cyan!15, opacity=0.7] (0,0) -- (2.5,2.5) -- (-2.5,2.5) -- cycle;
    \fill[quantumGold!15, opacity=0.7] (0,0) -- (2.5,-2.5) -- (-2.5,-2.5) -- cycle;
    
    % Labels
    \node at (0, 1.6) [font=\footnotesize\bfseries, align=center, color=darkNavy] {อนาคตสัมบูรณ์\\(Future Light Cone)};
    \node at (0, -1.6) [font=\footnotesize\bfseries, align=center, color=darkNavy] {อดีตสัมบูรณ์\\(Past Light Cone)};
    \node at (2.2, 0.8) [font=\scriptsize\color{purple}, align=center] {ย่านเหมือนปริภูมิ\\$\Delta s^2 < 0$};
    \node at (-2.2, 0.8) [font=\scriptsize\color{purple}, align=center] {ย่านเหมือนปริภูมิ\\$\Delta s^2 < 0$};
    
    % Light ray
    \node at (2.5, 2.8) [font=\scriptsize\bfseries\color{cyan}] {ลำแสง ($x = ct$)};
    \node at (-2.5, 2.8) [font=\scriptsize\bfseries\color{cyan}] {ลำแสง ($x = -ct$)};
    
    % Worldline
    \draw[red!80!black, line width=1.8pt, smooth] plot coordinates {(0,0) (0.3, 0.8) (0.2, 1.7) (0.6, 2.7)};
    \node at (0.8, 2.9) [font=\scriptsize\bfseries\color{red!80!black}] {เส้นทางโลก (Worldline)};
    
    % Event at origin
    \filldraw[darkNavy] (0,0) circle (2.5pt) node[below right, font=\footnotesize\bfseries] {ปัจจุบัน ($O$)};
\end{tikzpicture}
\caption{แผนภาพกรวยแสง (Light Cone) แสดงโครงสร้างการส่งผ่านผลเชิงเหตุและผล (Causality) ในปริภูมิ-เวลามินคอฟสกี}
\label{fig:light_cone}
\end{figure}

\newpage
\section{พลศาสตร์สัมพัทธภาพ โมเมนตัมและพลังงานรวม}
\index{โมเมนตัมสัมพัทธภาพ}
\index{พลังงานนิ่ง}
\index{สมมูลมวล-พลังงาน}

เมื่อวัตถุมีความเร็วเข้าใกล้อัตราเร็วแสง นิยามโมเมนตัมแบบนิวตัน $\mathbf{p} = m\mathbf{v}$ จำเป็นต้องได้รับการปรับปรุงตามทฤษฎีสัมพัทธภาพ เพื่อให้สอดคล้องกับกฎการอนุรักษ์โมเมนตัมในทุกกรอบเฉื่อย
\begin{equation}
\mathbf{p} = \gamma m_0 \mathbf{v} = \frac{m_0 \mathbf{v}}{\sqrt{1 - v^2/c^2}}
\end{equation}
โดยที่ $m_0$ คือมวลนิ่ง (Rest Mass) ของอนุภาค

พลังงานรวมสัมพัทธภาพ (Total Relativistic Energy, $E$) ประกอบด้วย \textbf{\mbox{พลังงานนิ่ง} (Rest Energy, $E_0$)} และ \textbf{พลังงานจลน์สัมพัทธภาพ ($K$)}
\begin{equation}
E = \gamma m_0 c^2 = K + m_0 c^2
\end{equation}
ดังนั้น พลังงานจลน์จึงมีค่าเท่ากับ
\begin{equation}
K = (\gamma - 1) m_0 c^2 = \left( \frac{1}{\sqrt{1 - v^2/c^2}} - 1 \right) m_0 c^2
\end{equation}
เมื่อ $v \ll c$ สามารถกระจายอนุกรมเทย์เลอร์ $\gamma \approx 1 + \frac{1}{2}\frac{v^2}{c^2}$ จะได้ $K \approx \frac{1}{2}m_0 v^2$ ซึ่งสอดคล้องกับกลศาสตร์คลาสสิกของนิวตันอย่างสมบูรณ์แบบ

ความสัมพันธ์พื้นฐานระหว่างพลังงานรวม โมเมนตัม และมวลนิ่ง กำหนดโดย
\begin{equation}
E^2 = (pc)^2 + (m_0 c^2)^2
\end{equation}
สำหรับอนุภาคที่ไร้มวลนิ่ง เช่น โฟตอน ($m_0 = 0$) สมการจะลดรูปเหลือ
\begin{equation}
E = pc
\end{equation}

\section{สถาปัตยกรรมการจำลองเสมือนจริง AR/XR 4D Spacetime Visualizer}
\index{WebXR}
\index{การจำลอง AR}

\begin{arbox}[การแปลงพิกัดปริภูมิ-เวลา 4 มิติ และกรวยแสงปฏิสัมพันธ์]
\textbf{วัตถุประสงค์การเรียนรู้ผ่าน AR}
\begin{enumerate}
    \item แสดงภาพกรวยแสง (Light Cone) แบบโปร่งใส 3 มิติ ลอยตัวในห้องเรียนจริงผ่าน WebXR
    \item ผู้เรียนสามารถเลื่อนแถบควบคุมความเร็ว $\beta = v/c$ ตั้งแต่ $0.00$ ถึง $0.99$ เพื่อสังเกตการเอียง (Shear / Hyperbolic Rotation) ของแกนพิกัด $x'$ และ $ct'$ เทียบกับแกน $x$ และ $ct$
    \item ตรวจสอบเส้นทางโลก (Worldline) ของอนุภาคความเร็วสูง และสังเกตการหดตัวของยานอวกาศเสมือนจริง
\end{enumerate}

\textbf{ข้อกำหนดทางเทคนิค (Technical Specification)}
\begin{itemize}
    \item \textbf{Shader Pipeline}\quad GLSL Fragment Shader คำนวณไฮเปอร์โบลา $\Delta s^2 = \text{const}$ แบบเรียลไทม์
    \item \textbf{Tracking}\quad MediaPipe Hands ควบคุมการหมุนมุมมองและ Pinch Gesture เพื่อเร่งความเร็วของยาน
    \item \textbf{Performance Target}\quad $60\text{ FPS}$ บนเบราว์เซอร์มือถือ รองรับ WebXR API
\end{itemize}
\end{arbox}

\section{การคำนวณเชิงตัวเลขด้วย Python การวิเคราะห์ตัวคูณลอเรนซ์}
\index{Python}
\index{NumPy}
\index{Matplotlib}

ต่อไปนี้เป็นสคริปต์ Python วิเคราะห์ค่าตัวคูณลอเรนซ์ $\gamma$ และพลังงานจลน์สัมพัทธภาพเทียบกับค่าคลาสสิก

\begin{pythonbox}[relativity\_kinematics.py -- การวิเคราะห์ตัวคูณลอเรนซ์และพลังงานจลน์]
import numpy as np
import matplotlib.pyplot as plt

# ค่าคงตัวมูลฐาน
c = 2.99792458e8     # อัตราเร็วแสง (m/s)
m_e = 9.1093837e-31   # มวลนิ่งของอิเล็กตรอน (kg)
eV = 1.602176634e-19  # จูลต่ออิเล็กตรอนโวลต์

# กำหนดช่วงความเร็วสัมพันธ์ beta = v/c จาก 0 ถึง 0.99
beta = np.linspace(0.001, 0.99, 1000)
gamma = 1.0 / np.sqrt(1.0 - beta**2)

# คำนวณพลังงานจลน์สัมพัทธภาพและพลังงานจลน์แบบคลาสสิก (หน่วย MeV)
K_rel_MeV = (gamma - 1.0) * m_e * c**2 / (1e6 * eV)
K_classical_MeV = 0.5 * m_e * (beta * c)**2 / (1e6 * eV)

# แสดงผลค่าตัวเลขที่สำคัญ
print("=== Relativistic Kinematics Calculation ===")
speeds = [0.1, 0.5, 0.8, 0.9, 0.95, 0.99, 0.999]
for v_rel in speeds:
    g = 1.0 / np.sqrt(1.0 - v_rel**2)
    k_rel = (g - 1.0) * 0.5109989  # พลังงานจลน์อิเล็กตรอน MeV
    k_cla = 0.5 * 0.5109989 * v_rel**2
    ratio = k_rel / k_cla
    print(f"beta = {v_rel:5.3f} | gamma = {g:7.3f} | K_rel = {k_rel:9.4f} MeV | Ratio = {ratio:6.2f}")
\end{pythonbox}

\section{ตัวอย่างการแก้โจทย์เชิงวิเคราะห์ตามระเบียบวิธี P-S-C-T}

\begin{psctexample}[การสลายตัวของมิวออนในบรรยากาศโลก (Muon Decay Paradox)]
\textbf{โจทย์ (Problem)} \\
มิวออน ($\mu$) ก่อตัวขึ้นในชั้นบรรยากาศสตราโตสเฟียร์ที่ระดับความสูง $H = 10\text{ km}$ เหนือพื้นโลก ด้วยความเร็ว $v = 0.995c$ โดยมิวออนมีอายุขัยเฉลี่ยขณะอยู่นิ่ง (Proper Lifetime) เท่ากับ $\tau_0 = 2.20\,\mu\text{s}$ 
\begin{enumerate}
    \item ตามการคำนวณแบบคลาสสิก มิวออนจะเคลื่อนที่ได้ระยะทางเท่าใดก่อนสลายตัว และสามารถมาถึงพื้นผิวโลกได้หรือไม่?
    \item จากมุมมองของผู้สังเกตบนพื้นโลก (Earth Frame) จงอธิบายว่าเหตุใดจึงตรวจพบมิวออนบนพื้นผิวโลกได้เป็นจำนวนมาก
    \item จากมุมมองของมิวออนเอง (Muon Frame) จงอธิบายว่ามิวออนเดินทางมาถึงพื้นโลกได้อย่างไร
\end{enumerate}

\textbf{ยุทธวิธี (Strategy)} \\
\begin{itemize}
    \item ในกรอบคลาสสิก ระยะทางคือ $d = v \tau_0$
    \item ในกรอบโลก เวลาของมิวออนยืดออก $\Delta t = \gamma \tau_0$ ดังนั้นระยะทางที่มิวออนเคลื่อนที่ได้คือ $d = v \Delta t$
    \item ในกรอบมิวออน ระยะทางความสูง $10\text{ km}$ หดสั้นลง $H' = H / \gamma$ แต่มิวออนมีเวลาสลายตัวเท่าเดิมคือ $\tau_0$
\end{itemize}

\textbf{การคำนวณ (Computation)} \\
คำนวณตัวคูณลอเรนซ์สำหรับความเร็ว $v = 0.995c$
\begin{equation}
\gamma = \frac{1}{\sqrt{1 - (0.995)^2}} = \frac{1}{\sqrt{1 - 0.990025}} = \frac{1}{\sqrt{0.009975}} \approx 10.0125
\end{equation}

\textit{1. การวิเคราะห์แบบคลาสสิก}
\begin{align}
d_{\text{classical}} &= v \tau_0 = (0.995 \times 3.00 \times 10^8\text{ m/s}) \times (2.20 \times 10^{-6}\text{ s}) \nonumber \\
&\approx 656.7\text{ m} \approx 0.66\text{ km}
\end{align}
เนื่องจาก $0.66\text{ km} \ll 10\text{ km}$ มิวออนเกือบทั้งหมดจึงควรสลายตัวไปหมดก่อนถึงพื้นโลก ซึ่งขัดแย้งกับข้อเท็จจริงจากการทดลอง

\textit{2. การวิเคราะห์ในกรอบโลก (Time Dilation)}
เวลาชีวิตของมิวออนที่สังเกตจากบนโลกจะยืดออกเป็น
\begin{equation}
\Delta t = \gamma \tau_0 = (10.0125) \times (2.20\,\mu\text{s}) \approx 22.03\,\mu\text{s}
\end{equation}
ระยะทางที่มิวออนเคลื่อนที่ได้ในกรอบโลกจึงเป็น
\begin{align}
d_{\text{Earth}} &= v \Delta t = (0.995 \times 3.00 \times 10^8\text{ m/s}) \times (22.03 \times 10^{-6}\text{ s}) \nonumber \\
&\approx 6576\text{ m} \approx 6.58\text{ km}
\end{align}
หากพิจารณาการกระจายทางสถิติของอายุขัย $N(t) = N_0 e^{-t/\Delta t}$ จะพบมิวออนจำนวนมากสามารถรอดชีวิตมาถึงสถานีตรวจวัดบนพื้นโลกได้จริง!

\textit{3. การวิเคราะห์ในกรอบมิวออน (Length Contraction)}
ในกรอบที่มิวออนอยู่นิ่ง โลกกำลังพุ่งเข้าหามันด้วยความเร็ว $0.995c$ ระยะทางความหนาของบรรยากาศจะหดสั้นลง
\begin{equation}
H' = \frac{H}{\gamma} = \frac{10\,000\text{ m}}{10.0125} \approx 998.75\text{ m} \approx 1.0\text{ km}
\end{equation}
เวลาที่โลกใช้ในการพุ่งผ่านมิวออนคือ
\begin{equation}
\Delta t' = \frac{H'}{v} = \frac{998.75\text{ m}}{0.995 \times 3.00 \times 10^8\text{ m/s}} \approx 3.35\,\mu\text{s}
\end{equation}
ซึ่งใกล้เคียงกับอายุขัย $\tau_0 = 2.20\,\mu\text{s}$ ทำให้มิวออนมีโอกาสชนพื้นโลกได้อย่างสมเหตุสมผล

\textbf{ข้อคิดทางฟิสิกส์ (Takeaway)} \\
ทั้งสองกรอบอ้างอิงให้ข้อสรุปทางกายภาพเดียวกันคือ "มิวออนตรวจพบบนพื้นผิวโลกได้" โดยกรอบโลกอธิบายด้วยปรากฏการณ์การยืดออกของเวลา (Time Dilation) ขณะที่กรอบมิวออนอธิบายด้วยการหดสั้นของความยาว (Length Contraction) นี่คือความสอดคล้องอย่างงดงามของทฤษฎีสัมพัทธภาพพิเศษ
\qedexam
\end{psctexample}

\begin{psctexample}[พลังงานขีดเริ่มเปลี่ยนในการสร้างอนุภาคไพออน (Threshold Energy for Pion Production)]
ลำอนุภาคโปรตอนพลังงานสูงจากเครื่องเร่งอนุภาคพุ่งชนเป้าโปรตอนที่อยู่นิ่งในห้องปฏิบัติการเพื่อสร้างไพออนที่เป็นกลางทางไฟฟ้าตามปฏิกิริยา
\begin{equation}
p + p \to p + p + \pi^0
\end{equation}
กำหนดให้มวลนิ่งของโปรตอน $m_p c^2 = 938.27\text{ MeV}$ และมวลนิ่งของไพออน $m_{\pi^0} c^2 = 134.98\text{ MeV}$
\begin{enumerate}
    \item จงหาสมการพลังงานรวมขีดเริ่มเปลี่ยน (Threshold Total Energy) $E_{\text{th}}$ และพลังงานจลน์ขีดเริ่มเปลี่ยน $K_{\text{th}}$ ของโปรตอนในกรอบห้องปฏิบัติการ
    \item คำนวณค่าตัวเลขของ $K_{\text{th}}$ ในหน่วย MeV
\end{enumerate}

\textbf{โจทย์ (Problem)} \\
คำนวณพลังงานจลน์ต่ำสุดที่ต้องให้กับโปรตอนตัวพุ่งชน เพื่อให้เกิดปฏิกิริยาสร้างอนุภาคใหม่ $\pi^0$ ขึ้นมาได้

\textbf{ยุทธวิธี (Strategy)} \\
ใช้ปริมาณคงตัวลอเรนซ์ของเวกเตอร์ 4-โมเมนตัมรวม (Mandelstam Invariant $s$)
\begin{equation}
s = (P_1 + P_2)^2 = E_{\text{tot, CM}}^2 / c^2
\end{equation}
\begin{itemize}
    \item ในกรอบห้องปฏิบัติการ (Lab Frame)\quad โปรตอนตัวที่สองอยู่นิ่ง $P_2 = (m_p c, \mathbf{0})$ และโปรตอนตัวแรกมี $P_1 = (E/c, \mathbf{p})$
    \begin{equation}
    s = (P_1 + P_2)^2 = 2 m_p^2 c^2 + 2 \left(\frac{E}{c}\right) (m_p c) = 2 m_p^2 c^2 + 2 m_p E
    \end{equation}
    \item ที่พลังงานขีดเริ่มเปลี่ยน ในกรอบจุดศูนย์กลางโมเมนตัม (CM Frame) อนุภาคผลลัพธ์ทั้งหมด ($p, p, \pi^0$) จะอยู่นิ่งสัมพัทธ์ซึ่งกันและกัน
    \begin{equation}
    s = (2 m_p + m_{\pi^0})^2 c^2
    \end{equation}
\end{itemize}

\textbf{การคำนวณ (Computation)} \\
\textit{1. เทียบปริมาณคงตัว $s$ ทั้งสองกรอบ}
\begin{align}
2 m_p^2 c^2 + 2 m_p E_{\text{th}} &= (2 m_p + m_{\pi^0})^2 c^2 \nonumber \\
&= 4 m_p^2 c^2 + 4 m_p m_{\pi^0} c^2 + m_{\pi^0}^2 c^2
\end{align}
แก้สมการหา $E_{\text{th}}$
\begin{equation}
E_{\text{th}} = m_p c^2 + 2 m_{\pi^0} c^2 + \frac{m_{\pi^0}^2 c^2}{2 m_p}
\end{equation}
พลังงานจลน์ขีดเริ่มเปลี่ยนคือ $K_{\text{th}} = E_{\text{th}} - m_p c^2$
\begin{equation}
K_{\text{th}} = 2 m_{\pi^0} c^2 + \frac{(m_{\pi^0} c^2)^2}{2 m_p c^2}
\end{equation}

\textit{2. แทนค่าตัวเลข}
\begin{align}
K_{\text{th}} &= 2 (134.98\text{ MeV}) + \frac{(134.98\text{ MeV})^2}{2 (938.27\text{ MeV})} \nonumber \\
&= 269.96\text{ MeV} + \frac{18219.60}{1876.54}\text{ MeV} \nonumber \\
&= 269.96\text{ MeV} + 9.71\text{ MeV} \approx 279.67\text{ MeV}
\end{align}

\textbf{ข้อคิดทางฟิสิกส์ (Takeaway)} \\
มวลของอนุภาค $\pi^0$ ที่ต้องการสร้างคือ $134.98\text{ MeV}/c^2$ แต่ต้องใช้พลังงานจลน์ของโปรตอนตัวพุ่งชนถึง $279.67\text{ MeV}$ (มากกว่าสองเท่าของมวลที่สร้างขึ้น) เนื่องจากในกรอบห้องปฏิบัติการ อนุภาคทั้งหมดต้องเคลื่อนที่ไปข้างหน้าเพื่อรักษาโมเมนตัมรวมของระบบไว้ พลังงานจลน์ส่วนใหญ่จึงสูญเสียไปกับการเคลื่อนที่ของจุดศูนย์กลางมวล นี่คือเหตุผลที่ฟิสิกส์พลังงานสูงในปัจจุบันหันมาใช้เครื่องชนอนุภาคแบบลำอนุภาคปะทะกัน (Colliders) เช่น LHC แทนเครื่องเร่งอนุภาคแบบชนเป้านิ่ง
\qedexam
\end{psctexample}

\begin{psctexample}[การปรับค่าพิกัดการแปลงพื้นที่สำหรับ AR Tracking ให้เรียบเนียน (Smooth Interpolation / LERP)]
\textbf{โจทย์ (Problem)} \\
ในการจำลองทฤษฎีสัมพัทธภาพพิเศษด้วยระบบ AR เมื่อเซนเซอร์ตรวจจับตำแหน่ง (เช่น WebXR หรือ MediaPipe Hands) ส่งพิกัดตำแหน่งของกรอบอ้างอิงเฉื่อย $(\mathbf{x}_t, \mathbf{y}_t, \mathbf{z}_t)$ เข้ามาประมวลผลทุกเฟรม มักพบปัญหาการสั่นไหวเล็กน้อยของข้อมูลจากเซนเซอร์ (Jittering / Sensor Noise) ส่งผลให้กรอบอวกาศ-เวลา 4 มิติ และกรวยแสงเสมือนสั่นกระตุกจนสูญเสียความสมจริง จงออกแบบตัวกรองทางคณิตศาสตร์ที่คำนวณได้อย่างรวดเร็วในรอบเฟรม 60 FPS เพื่อให้การเคลื่อนที่ราบรื่น

\textbf{ยุทธวิธี (Strategy)} \\
ใช้ตัวกรองแบบชะลอชี้กำลังหรือการประมาณค่าในช่วงเชิงเส้น (Linear Interpolation -- LERP) ระหว่างตำแหน่งปัจจุบัน $\mathbf{P}_{\text{curr}}$ และตำแหน่งเป้าหมาย $\mathbf{P}_{\text{target}}$
\begin{equation}
\mathbf{P}(t + \Delta t) = \mathbf{P}(t) + \alpha \left( \mathbf{P}_{\text{target}} - \mathbf{P}(t) \right) = (1 - \alpha)\mathbf{P}(t) + \alpha \mathbf{P}_{\text{target}}
\end{equation}
โดยที่ $\alpha \in (0, 1]$ คือแฟกเตอร์การผ่อนคลาย (Smoothing Factor)

\textbf{การคำนวณ (Computation)} \\
หากต้องการให้การตอบสนองมีความหน่วง 95\% เข้าสู่เป้าหมายภายในเวลา $T = 0.2\text{ วินาที}$ ที่ความถี่เฟรม $60\text{ Hz}$ ($\Delta t = 1/60\text{ s}$ หรือมีจำนวน 12 เฟรม)
\begin{equation}
(1 - \alpha)^{12} \le 0.05 \implies 1 - \alpha = (0.05)^{1/12} \approx 0.779 \implies \alpha \approx 0.22
\end{equation}
ดังนั้น การเขียนโค้ด LERP ในลูปแอนิเมชันของ Three.js คือ
\begin{equation}
\mathbf{P}_{\text{curr}}.\text{lerp}(\mathbf{P}_{\text{target}}, 0.22)
\end{equation}

\textbf{ข้อคิดทางฟิสิกส์และวิศวกรรม (Takeaway)} \\
การคำนวณ LERP แบบเวกเตอร์ใช้การบวกและคูณสเกลาร์เพียง 3 มิติ ซึ่งใช้เวลาประมวลผลน้อยกว่า $0.001\text{ ms}$ บน CPU จึงไม่กระทบต่อ Frame Budget ของเบราว์เซอร์ แต่สามารถกำจัดความสั่นไหวของเซนเซอร์ AR ในการจำลองการเปลี่ยนกรอบอ้างอิงเฉื่อยได้อย่างมีประสิทธิภาพสูงสุด
\qedexam
\end{psctexample}

\begin{conceptcheck}[คำถามทบทวนแนวคิด]
\begin{enumerate}
    \item ยานอวกาศลำหนึ่งมีความเร็วสัมพัทธ์ $0.8c$ เมื่อเทียบกับสถานีอวกาศ หากนักบินอวกาศบนยานยิงแสงเลเซอร์ไปข้างหน้า ผู้สังเกตบนสถานีอวกาศจะวัดอัตราเร็วของลำแสงเลเซอร์ได้เท่าใด?
    \item มีเงื่อนไขใดหรือไม่ที่ทำให้เหตุการณ์สองเหตุการณ์เกิดขึ้นพร้อมกันในทุกกรอบอ้างอิงเฉื่อย? จงอธิบายโดยใช้สมการการแปลงลอเรนซ์
\end{enumerate}
\end{conceptcheck}

\section{สรุปสาระสำคัญประจำบท}

\begin{table}[H]
\centering
\caption{สรุปความสัมพันธ์พื้นฐานในทฤษฎีสัมพัทธภาพพิเศษ}
\label{tab:relativity_summary}
\small
\begin{tabularx}{\textwidth}{l >{\centering\arraybackslash}p{4.2cm} Y}
\toprule
\textbf{หัวข้อ} & \textbf{สมการหลัก} & \textbf{ความหมายเชิงกายภาพ} \\
\midrule
ตัวคูณลอเรนซ์ & $\gamma = \frac{1}{\sqrt{1 - v^2/c^2}}$ & สเกลาร์มาตรวัดความเบี่ยงเบนจากกลศาสตร์นิวตัน ($\gamma \ge 1$) \\
การยืดของเวลา & $\Delta t = \gamma \Delta t_0$ & เวลาในกรอบที่วัตถุเคลื่อนที่เดินช้ากว่าเวลาเฉพาะ \\
การหดสั้นของความยาว & $L = L_0 / \gamma$ & ความยาวในทิศทางการเคลื่อนที่จะหดสั้นลงเมื่อสังเกตจากกรอบอื่น \\
ช่วงปริภูมิ-เวลา & $\Delta s^2 = c^2\Delta t^2 - |\Delta\mathbf{r}|^2$ & ปริมาณคงตัวลอเรนซ์ จำแนกย่านแสง เวลา และปริภูมิ \\
โมเมนตัมสัมพัทธภาพ & $\mathbf{p} = \gamma m_0 \mathbf{v}$ & กฎอนุรักษ์โมเมนตัมคงรูปในทุกกรอบอ้างอิงเฉื่อย \\
สมมูลมวล-พลังงาน & $E^2 = (pc)^2 + (m_0 c^2)^2$ & พลังงานรวมของอนุภาคที่มีมวลและอนุภาคไร้มวล \\
\bottomrule
\end{tabularx}
\end{table}

\section{แบบฝึกหัดท้ายบท}
\begin{enumerate}
    \item โปรตอนถูกเร่งในเครื่องเร่งอนุภาคจนมีพลังงานจลน์เป็น 3 เท่าของพลังงานนิ่ง จงคำนวณ
    \begin{enumerate}
        \item ตัวคูณลอเรนซ์ $\gamma$ ของโปรตอน
        \item อัตราเร็วของโปรตอนในรูปสัดส่วนของ $c$
        \item โมเมนตัมสัมพัทธภาพของโปรตอนในหน่วย $\text{MeV}/c$ (กำหนด $m_p c^2 \approx 938.27\text{ MeV}$)
    \end{enumerate}
    \item ยานอวกาศยาว $100\text{ m}$ เคลื่อนที่ผ่านสถานีอวกาศด้วยความเร็ว $0.6c$ ผู้สังเกตบนสถานีจะวัดความยาวของยานลำนี้ได้เท่าใด?
    \item จงพิสูจน์ว่าช่วงปริภูมิเวลา $\Delta s^2 = c^2 \Delta t^2 - \Delta x^2$ มีค่าเท่าเดิมภายใต้การแปลงลอเรนซ์ ($c^2 \Delta t'^2 - \Delta x'^2 = c^2 \Delta t^2 - \Delta x^2$)
    \item ปรากฏการณ์ด็อปเปลอร์เชิงสัมพัทธภาพ (Relativistic Doppler Effect) สำหรับแสง กาแล็กซีอันไกลโพ้นแห่งหนึ่งมีสเปกตรัมเส้นไฮโดรเจน $H_\alpha$ ($\lambda_0 = 656.3\text{ nm}$) ปรากฏที่ความยาวคลื่นที่วัดได้บนโลก $\lambda = 800.0\text{ nm}$ (เกิดการเลื่อนทางแดง Redshift) จงคำนวณอัตราเร็วในการถอยห่าง $v$ ของกาแล็กซีนี้ในรูปของสัดส่วนอัตราเร็วแสง $c$
    \item ยานอวกาศ $A$ และยานอวกาศ $B$ เคลื่อนที่ออกจากสถานีอวกาศในทิศทางตรงกันข้ามกันด้วยอัตราเร็ว $0.70c$ และ $0.80c$ ตามลำดับ เมื่อเทียบกับสถานีอวกาศ จงคำนวณ
    \begin{enumerate}
        \item อัตราเร็วของยาน $B$ เมื่อสังเกตจากผู้สังเกตบนยาน $A$ ตามทฤษฎีสัมพัทธภาพพิเศษ
        \item เปรียบเทียบผลลัพธ์ที่ได้กับการคำนวณตามจลนศาสตร์แบบคลาสสิกของกาลิเลโอ
    \end{enumerate}
\end{enumerate}
